% !TeX spellcheck = es_ES

%------------Basic document settings--------------%
\documentclass[12pt, openany]{article}
\usepackage[a4paper,width = 150mm, top=25mm, bottom= 25mm, footskip=35pt]{geometry}
\usepackage[utf8]{inputenc}
\usepackage[spanish]{babel}

%---------Graphics paths and labels------------%
\usepackage{graphicx}
\graphicspath{{Images/}}
\usepackage[labelfont=bf]{caption}
\usepackage{subcaption}
\usepackage{titlesec}

%-Miscellaneous packages for different purposes-%
\usepackage{chngcntr}
\usepackage{setspace}
\usepackage{fancyhdr}
\usepackage{url}
\usepackage{lastpage} 
\usepackage[bottom]{footmisc}
\usepackage{pdflscape}
\usepackage{ragged2e}
\usepackage{enumerate}
\usepackage{enumitem} 
\usepackage{todonotes}
\usepackage{multirow}
\usepackage{bm}
\usepackage{afterpage}
\usepackage{tabularx}
\usepackage{float}

%-------------Math and Physics packages-----------%
\usepackage{amssymb}
\usepackage{amsthm}
\usepackage{amsmath}
\usepackage{amsfonts}
\usepackage{mhchem}
\usepackage{physics}
\usepackage{mathtools}
\usepackage{empheq}

%---------Sets up different colors----------------%

\usepackage{hyperref}
\usepackage{xcolor}
\usepackage{layouts}
\hypersetup{
	colorlinks=true,
	linkcolor=blue,
	filecolor=magenta,      
	urlcolor=blue,
	citecolor=blue,
}
\usepackage{listings}
\usepackage{color} %red, green, blue, yellow, cyan, magenta, black, white
\definecolor{mygreen}{RGB}{28,172,0} % color values Red, Green, Blue
\definecolor{mylilas}{RGB}{170,55,241}
%\usepackage{tabu}

%---------Set specific bibliography--------------%

\usepackage{natbib}
\bibliographystyle{plainnat}
\setcitestyle{authoryear,open={(},close={)}}

%--------Useful for Mathematical Texts-----------%

\newtheorem{theorem}{Theorem}[section]
\newtheorem{corollary}{Corollary}[theorem]
\newtheorem{lemma}[theorem]{Lemma}
\newtheorem{definition}{Definition}[section]
\newtheorem{prop}{Proposition}
\newtheorem{property}{Property}
\theoremstyle{definition}
\newtheorem{remark}{Remark}[section]
\newtheorem{example}{Example}[section]

%--------Definition of divergente operator-------%
\AtBeginDocument{
	\let\div\relax
	\DeclareMathOperator{\div}{div}
}
\DeclareMathOperator{\ran}{ran}
\newcommand*\df{\mathop{}\!\mathrm{d}}
\newcommand*\Df[1]{\mathop{}\!\mathrm{d^#1}}

%---------Header and footer definitions----------%

% Define Header
\fancyhf{}
\fancyhead[L]{\hspace{0.3cm}{\leftmark}} 

% Define Footer
\fancyfoot[L]{\hspace{0.3cm} Álvaro Zurdo Navajo}
\fancyfoot[C]{\thepage}
\fancyfoot[R]{Trabajo 1 - MCPS \hspace{0.3cm}} 

%Define width of horizontal line in the header
\renewcommand{\headrulewidth}{0.4pt}
\setlength{\headsep}{0.2in}

%Define width of horizontal line in the footer
\renewcommand{\footrulewidth}{0.4pt}
\setlength{\headheight}{25pt}

%--Establish when figures and tables begin to count--%

\counterwithin{figure}{section}
\counterwithin{table}{section}
\numberwithin{equation}{section}

%--Pasar algunas palabras a español--%
%\renewcommand{\listtablename}{Índice de tablas}
\addto\captionsspanish{%
  \renewcommand{\listtablename}{Lista de tablas}%
  \renewcommand{\listfigurename}{Lista de figuras}%
  \renewcommand{\tablename}{Tabla}%
  \renewcommand{\lstlistingname}{Código}
}


%----------Use of nomenclature section-----------%
\usepackage{enumitem}
\usepackage{nomencl}
\usepackage{etoolbox}
\usepackage{acronym}
%\renewcommand\nomgroup[1]{%
%	\item[\bfseries
%	\ifstrequal{#1}{A}{Chapter 2}{%
%	\ifstrequal{#1}{B}{Chapter 3}{%
%	\ifstrequal{#1}{C}{Chapter 4}{}}}%
%	]}
\renewcommand\nomgroup[1]{%
	\item[\bfseries
	\ifstrequal{#1}{A}{Notation}{%
	\ifstrequal{#1}{B}{Symbols}{}}%
	]}
\makenomenclature

%----------Code formating-----------%

\definecolor{codegreen}{rgb}{0,0.6,0}
\definecolor{codegray}{rgb}{0.5,0.5,0.5}
\definecolor{codepurple}{rgb}{0.58,0,0.82}
\definecolor{backcolour}{rgb}{0.95,0.95,0.92}

\lstdefinestyle{mystyle}{
    backgroundcolor=\color{backcolour},
    commentstyle=\color{codegreen},
    keywordstyle=\color{magenta},
    numberstyle=\tiny\color{codegray},
    stringstyle=\color{codepurple},
    basicstyle=\footnotesize,
    breakatwhitespace=false,
    breaklines=true,
    captionpos=b,
    keepspaces=true,
    numbers=left,
    numbersep=5pt,
    showspaces=false,
    showstringspaces=false,
    showtabs=false,
    tabsize=2
}

\lstset{style=mystyle}

\begin{document}

\pagenumbering{roman}
\onehalfspacing

\begin{center}
\thispagestyle{empty}


\begin{figure}
	\begin{subfigure}{0.5\textwidth}
		\includegraphics[width=\linewidth]{Images/Portada/Logo_UPM.png}\\
	\end{subfigure}
	\begin{subfigure}{0.5\textwidth}
		\includegraphics[width=\linewidth]{Images/Portada/Logo_ETSIAE.png}\\
	\end{subfigure}
%	\begin{subfigure}{0.32\textwidth}
%		\includegraphics[width=\linewidth]{Images/Portada/Logo_MUSE.png}\\
%	\end{subfigure}
\end{figure}

\begin{figure}
\centering
\includegraphics[width=0.75\linewidth]{Images/Portada/Logo_MUSE.png}
\end{figure}

\vspace*{3cm} % Este comando no funciona, comprobar por qué
{\huge \textbf{Título}}\\
\vspace{0.75cm}
\textbf{Asignatura} \\
\begin{center}
    \large{\textbf{\today}} \\
\end{center}

\vspace{5cm} % 6.5

\begin{center}
        Álvaro Zurdo Navajo
\end{center}

\end{center}
\clearpage

\doublespacing	
\tableofcontents
\pagestyle{plain}
\clearpage

%\addcontentsline{toc}{section}{\listfigurename}
%\listoffigures
%\clearpage
%\thispagestyle{plain}
%
%\addcontentsline{toc}{section}{\listtablename}
%\listoftables
%\clearpage

\pagestyle{fancy}
\pagenumbering{arabic}
\setcounter{page}{1}	
\onehalfspacing		
\section{Introducción}
	\label{chap:intro}
	% !TeX spellcheck = en_US


\clearpage
	
\section{Desarrollo}
	\label{chap:1}
	\input{Chapters/Chapter1}
\clearpage

\section{Código}
	\label{chap:2}
	% !TeX spellcheck = en_US


\clearpage

\section{Resultados}
	\label{chap:resultados}
	\input{Chapters/Resultados}

\end{document}}
		